\FloatBarrier
% ============================================================
% 05-experimental-results.tex
% ============================================================

\section{Experimental Results}
\label{sec:exp}
\suppressfloats[t]

This section presents the experimental dataset, implementation environment,
evaluation metrics, experimental scenarios, and discussion of the proposed
CTU-Chat framework.

%------------------------------------------------
\subsection{Experimental Dataset}

A benchmark dataset was constructed to evaluate the proposed CTU-Chat
framework on university student-service question answering. The dataset was
developed from official Can Tho University documents, including academic
regulations, training programs, tuition policies, scholarships, student loan
policies, social support, dormitory regulations, and administrative procedures.

Instead of manually composing every question, candidate questions were first
generated using Claude based on the collected documents. The generated dataset
then underwent a two-stage validation process. First, a validation notebook was
used to verify the availability of supporting documents and evidence references.
Second, all questions were manually reviewed to remove duplicated or ambiguous
samples and to ensure that each question had a complete ground-truth answer.
A final validation step was performed before constructing the benchmark.

The final benchmark contains 150 representative questions selected according to
three criteria: (1) coverage of all major administrative categories,
(2) high-frequency real-world student inquiries, and (3) complete
ground-truth evidence for every sample.

\begin{table}[!htbp]
\centering
\caption{Distribution of the benchmark dataset.}
\label{tab:dataset}
\begin{tabular}{lc}
\toprule
Category & Questions \\
\midrule
Actual Tuition & 60 \\
Academic Rules & 18 \\
Scholarship & 14 \\
Student Loan & 12 \\
Other & 12 \\
Social Support & 11 \\
Academic Program & 9 \\
Exemption Policy & 9 \\
Exemption Basis & 5 \\
\midrule
\textbf{Total} & \textbf{150} \\
\bottomrule
\end{tabular}
\end{table}

%------------------------------------------------
\subsection{Experimental Tools}

The proposed system was implemented using a hybrid retrieval architecture
consisting of vector, graph, and relational databases. Dense embeddings were
generated using the Vietnamese bi-encoder used by the deployed index, and a
BGE reranker was used in the upgraded retrieval pipeline. LangGraph was used
for multi-agent orchestration. Retrieval metrics were computed from gold
document identifiers; answer relevancy and correctness were scored by a
Gemini-based LLM judge.

\begin{table}[!htbp]
\centering
\caption{Experimental implementation environment.}
\label{tab:tools}
\begin{tabular}{ll}
\toprule
Component & Technology \\
\midrule
Programming language & Python 3.12 \\
Vector database & Qdrant \\
Graph database & Neo4j \\
Metadata database & PostgreSQL \\
Embedding model & \texttt{vietnamese-bi-encoder} \\
Cross-encoder reranker & \texttt{bge-reranker-v2-m3} \\
Generation and judge LLM & Gemini 2.5 Flash Lite \\
Agent framework & LangGraph \\
Evaluation & Gold-source metrics + LLM judge \\
\bottomrule
\end{tabular}
\end{table}

%------------------------------------------------
\subsection{Evaluation Metrics}

Two groups of evaluation metrics were adopted in this study.

\textbf{Retrieval metrics} evaluate the quality of evidence retrieval,
including Hit@1 (H@1), Hit@3 (H@3), Precision@5 (P@5), Recall@5, and
Mean Reciprocal Rank (MRR@10).

\textbf{End-to-end QA metrics} follow the terminology of Answer Relevancy
(AR), Context Recall (CR), Context Precision (CP), and Answer Correctness
(AC). In the reported run, CR and CP were computed deterministically from
gold document identifiers, whereas AR and AC were assigned by a Gemini 2.5
Flash Lite judge. These are therefore RAGAS-style metrics, not scores produced
by the RAGAS software package. All metrics range from 0 to 1, where higher
values indicate better performance.

%------------------------------------------------
\subsection{Experimental Scenarios}

Two groups of experiments were conducted using the same 150-question benchmark.
The first evaluates the retrieval module (E1--E5), while the second evaluates
the complete question answering pipeline (T1--T7).

%-----------------------------
\subsubsection{Scenario A: Retrieval Performance}

Five cumulative retrieval configurations were evaluated. E1 uses BM25, E2
uses dense retrieval, and E3 combines both retrievers through RRF. E4 adds the
\texttt{bge-reranker-v2-m3} cross-encoder, while E5 represents the complete
system with governed retrieval lanes, graph-based candidate expansion, and
cross-encoder reranking.

Table~\ref{tab:retrieval} shows a clear improvement as the reranker and full
system components are introduced. E4 raises H@1 from 0.6000 to 0.8000 and
MRR@10 from 0.7006 to 0.8556 relative to E3. E5 achieves the strongest overall
retrieval result, including perfect H@3 and Recall@5 and an MRR@10 of 0.9222.

\begin{table}[!ht]
\centering
\caption{Retrieval evaluation results of E1--E5 on 150 questions.}
\label{tab:retrieval}
\begin{tabular}{lccccc}
\toprule
Config. & H@1 & H@3 & P@5 & R@5 & MRR \\
\midrule
E1 & 0.7333 & 0.8667 & 0.2133 & 0.9000 & 0.8133 \\
E2 & 0.5333 & 0.6000 & 0.1600 & 0.6667 & 0.5833 \\
E3 & 0.6000 & 0.8000 & 0.2000 & 0.8333 & 0.7006 \\
E4 & 0.8000 & 0.9333 & 0.2267 & 0.9333 & 0.8556 \\
\textbf{E5} & \textbf{0.8667} & \textbf{1.0000} & \textbf{0.2400} & \textbf{1.0000} & \textbf{0.9222} \\
\bottomrule
\end{tabular}
\end{table}

%-----------------------------
\subsubsection{Scenario B: End-to-End Question Answering}

The second experiment evaluates seven QA configurations (T1--T7) using the
metric procedure described above. The logged run contains 150 cases for each
configuration, for a total of 1,050 generated answers and 1,050 judge records.

Table~\ref{tab:ragas} shows that T4 achieves strong context quality, with
Context Recall of 0.745 and Context Precision of 0.749. In contrast, removing
the governance filter in T7 reduces Answer Correctness to 0.398, the lowest
score in the experiment. T5 and T6 remain competitive on answer correctness,
indicating that high-quality retrieval and answer generation must be optimized
jointly.

\begin{table}[!ht]
\centering
\caption{End-to-end QA results of T1--T7 on 150 questions.}
\label{tab:ragas}
\begin{tabular}{lcccc}
\toprule
Config. & AR & CR & CP & AC \\
\midrule
T1 & 0.561 & 0.745 & 0.684 & 0.518 \\
T2 & 0.455 & 0.533 & 0.421 & 0.395 \\
T3 & 0.578 & 0.682 & 0.550 & 0.506 \\
\textbf{T4} & \textbf{0.551} & \textbf{0.745} & \textbf{0.749} & \textbf{0.500} \\
T5 & 0.575 & 0.682 & 0.550 & 0.525 \\
T6 & 0.561 & 0.745 & 0.750 & 0.508 \\
T7 & 0.479 & 0.533 & 0.421 & \textbf{0.398} \\
\bottomrule
\end{tabular}
\end{table}

%------------------------------------------------
\subsection{Discussion}

The retrieval results demonstrate the complementary roles of lexical matching,
semantic retrieval, and learned reranking. BM25 remains effective for exact
cohort identifiers and administrative terminology, whereas the cross-encoder
substantially improves the ranking of hybrid candidates. Governed lanes and
graph expansion further improve evidence coverage, allowing E5 to retrieve all
required sources within the top five for this benchmark.

The end-to-end results reinforce the importance of governed retrieval. T4
maintains high Context Recall and Context Precision, while T7 records the
lowest Answer Correctness. At the same time, the competitive AC values of T5
and T6 show that context quality alone does not guarantee the best generated
answer; generation prompts and answer calibration remain important optimization
targets.

The benchmark is drawn from a single institution and contains a large tuition
subset. Future work will extend the evaluation to balanced bilingual queries,
additional institutions, and repeated or human-assisted answer assessment.
